\textit{Design representation and boundary conditions.}
The support \(\mathcal Z_n\) is finite and invariant under global sign reversal. For a blinded law \(\pi\in\Pi_n\), define
\[
\alpha_s=\pi(s)+\pi(-s),\qquad s\in\mathscr S.
\]
Then \(\alpha_s\ge0\), \(\sum_s\alpha_s=1\), and \((-s)(-s)^\top=ss^\top\) gives
\[
Q_\pi=\mathbb E_\pi[ZZ^\top]
=\sum_{s\in\mathscr S}\alpha_sss^\top.
\tag{1}
\]
Conversely, every \(\alpha\in\Delta_m\) defines the blinded law in \((\mathrm D)\), and its second moment is exactly \((1)\). Thus the simplex of sign-pair masses parametrizes every second moment relevant to the criterion. Arbitrary dependence across blocks is allowed, as required by \(\Pi_n\). Zero masses are admissible: the argument uses no inverse assignment probability, so no additional overlap condition is needed.

Each \(G_j\) is nonzero and positive semidefinite, hence \(r_j=\operatorname{rank}(G_j)\ge1\), making division by \(r_j\) legitimate. Moore--Penrose inverses remain well defined when \(G_j\) is singular, and no commutativity assumption is imposed.

\textit{Generalized spectral risk and the causal variance bound.}
Fix \(j\). By the spectral decomposition of \(G_j\), every \(x\in\operatorname{range}(G_j)\) has a unique representation \(x=G_j^{1/2}u\) with \(u\in\operatorname{range}(G_j)\), and
\[
x^\top G_j^\dagger x=\lVert u\rVert_2^2.
\tag{2}
\]
Combining \((1)\), \((2)\), and the Rayleigh--Ritz identity yields
\[
\begin{aligned}
\sup_{x\in\mathcal M_j^{\mathrm{gen}}}
\frac1{n^2}x^\top Q_\pi x
&=\frac1{n^2}
\sup_{\substack{u\in\operatorname{range}(G_j)\\\lVert u\rVert_2\le1}}
u^\top G_j^{1/2}Q_\pi G_j^{1/2}u\\
&=\frac1{n^2}\lambda_{\max}
\{G_j^{1/2}Q_\pi G_j^{1/2}\}
=\frac1{n^2}\lambda_{\max}\{A_j(\alpha)\}.
\end{aligned}
\tag{3}
\]
The middle equality is valid because the displayed positive-semidefinite matrix has range contained in \(\operatorname{range}(G_j)\); the last equality follows from \((1)\) and the definition of \(A_{js}\). This argument is exact for noncommuting matrices.

Under \(\texttt{ass:fixed-potential-outcomes}\), \(\texttt{ass:no-interference}\), and \(\texttt{ass:assignment-only-randomness}\), \(\texttt{lem:variance-and-spectral-risk}\) gives
\[
\operatorname{Var}_\pi(\widehat\tau)
=n^{-2}\mu^\top Q_\pi\mu.
\tag{4}
\]
Therefore \(\mu\in\mathcal M_j^{\mathrm{gen}}\) and \((3)\)--\((4)\) imply the variance inequality in the statement. This calculation keeps the fixed causal estimand \(\tau_N\) distinct from the design-selection functional \(\lambda_{\mathrm{gen}}^\star\).

The simplex is compact and the largest-eigenvalue map is continuous, so each minimum defining \(\rho_j\) is attained. Equivalently,
\[
\rho_j=\min_{\alpha\in\Delta_m,\ u\in\mathbb R}
\left\{u:A_j(\alpha)\preceq n^2uI_n\right\}.
\tag{5}
\]
Thus \(\beta_j\), \(\rho_j\), and \(\delta_j\) are finite computable functions of the known menu and baseline law. The hypothesis \(\delta_j>0\) is exactly the positivity condition required to divide by the gain denominator. If it fails, the frozen gain fraction is undefined unless a separate degeneracy convention is supplied.

\textit{Exact primal program.}
For a real \(t\), the retained-gain inequality for class \(j\) is, by \((3)\),
\[
\frac{\beta_j-n^{-2}\lambda_{\max}\{A_j(\alpha)\}}{\delta_j}\ge t
\quad\Longleftrightarrow\quad
A_j(\alpha)\preceq n^2(\beta_j-\delta_jt)I_n.
\tag{6}
\]
Taking all class constraints simultaneously and maximizing over the common mass vector proves \((\mathrm P)\). The baseline mass vector \(\alpha^0\) is feasible at \(t=0\). Moreover, the definition of \(\rho_j\) implies that every class gain is at most one. Hence
\[
0\le\lambda_{\mathrm{gen}}^\star\le1.
\tag{7}
\]
For each \(\alpha\), the greatest feasible \(t\) is a continuous function of \(\alpha\); compactness of \(\Delta_m\) therefore gives a primal optimizer.

\textit{Trace relaxation and exactness.}
The range of \(A_j(\alpha)\) is contained in \(\operatorname{range}(G_j)\). Applying \(\texttt{lem:finite-trace-isotropy}\) gives
\[
\lambda_{\max}\{A_j(\alpha)\}
\ge \frac{\operatorname{tr}\{A_j(\alpha)\}}{r_j}.
\tag{8}
\]
Replacing the left side of \((6)\) by the lower bound in \((8)\) produces the finite linear program \((\mathrm T)\), and relaxation gives
\[
\lambda_{\mathrm{gen}}^\star\le\overline\lambda.
\tag{9}
\]

If equality holds in \((9)\), a primal optimizer at \(t=\overline\lambda\) satisfies exactly the spectral inequalities in \((\mathrm E)\). Conversely, an \(\alpha\) satisfying \((\mathrm E)\) is feasible for \((\mathrm P)\) at \(t=\overline\lambda\), so \((9)\) must be an equality. This proves the first exactness criterion.

Any \(\alpha\) satisfying \((\mathrm E)\) also satisfies all trace constraints at \(t=\overline\lambda\), and hence lies on the optimal face of \((\mathrm T)\). If class \(j\) is trace-active there, then
\[
\frac{\operatorname{tr}\{A_j(\alpha)\}}{r_j}
=n^2(\beta_j-\delta_j\overline\lambda).
\tag{10}
\]
Equations \((\mathrm E)\), \((8)\), and \((10)\) force equality in the trace bound. The equality clause of \(\texttt{lem:finite-trace-isotropy}\) is precisely
\[
A_j(\alpha)
=\frac{\operatorname{tr}\{A_j(\alpha)\}}{r_j}
P_{\operatorname{range}(G_j)},
\]
which is \((\mathrm I)\). A trace-slack class need only satisfy its spectral threshold in \((\mathrm E)\). Conversely, isotropy gives \((\mathrm E)\) for every trace-active class, and the stipulated threshold gives it for every trace-slack class. At least one trace constraint is active at an optimum because every \(\delta_j>0\); otherwise \(t\) could be increased. These conditions are therefore necessary and sufficient, and simultaneous isotropy is not demanded of slack classes.

\textit{Sharp gap, attaining design, and dual certificate.}
If the exactness condition fails, the proved equivalence makes \((9)\) strict. Thus
\[
\Gamma_n=\overline\lambda-\lambda_{\mathrm{gen}}^\star>0
\]
is the sharp additive trace-relaxation gap. It is computable as the difference between the two displayed finite optimization problems, not an unspecified extremum. If \((\alpha^\star,t^\star)\) solves \((\mathrm P)\), the blinded law \((\mathrm D)\) belongs to \(\Pi_n\), has the second moment in \((1)\), and, by \((3)\)--\((6)\), attains \(t^\star=\lambda_{\mathrm{gen}}^\star\). It therefore pays exactly \(\Gamma_n\), with no rounding loss.

It remains to verify the dual formula. Put
\[
F_j(\alpha,t)=n^2(\beta_j-\delta_jt)I_n-A_j(\alpha).
\]
For any \(X_j\succeq0\) satisfying
\[
n^2\sum_j\delta_j\operatorname{tr}(X_j)=1,
\tag{11}
\]
every primal-feasible \((\alpha,t)\) obeys
\[
\begin{aligned}
t
&\le n^2\sum_j\beta_j\operatorname{tr}(X_j)
-\sum_{s\in\mathscr S}\alpha_s
  \sum_j\operatorname{tr}(X_jA_{js})\\
&\le n^2\sum_j\beta_j\operatorname{tr}(X_j)
-\min_{s\in\mathscr S}\sum_j\operatorname{tr}(X_jA_{js}).
\end{aligned}
\tag{12}
\]
Introducing \(\gamma\le\sum_j\operatorname{tr}(X_jA_{js})\) for every \(s\) gives weak duality with the program in \((\mathrm C)\).

For the reverse inequality, choose \(\alpha\) in the relative interior of \(\Delta_m\), so all \(\alpha_s>0\), and take \(t\) sufficiently negative. Since every \(\delta_j>0\), all matrices \(F_j(\alpha,t)\) are then positive definite. Thus the finite-dimensional conic program is strictly feasible relative to its simplex equality constraint. To see the resulting multiplier form directly, separate a point with objective strictly above the primal optimum from the convex set of attainable tuples consisting of the matrix slacks, the simplex equality residual, and the objective hypograph. Strict feasibility makes the matrix-separating components positive semidefinite and rules out a zero objective multiplier. Normalize that multiplier to one. The coefficient of the free variable \(t\) then gives exactly \((11)\); minimizing the remaining affine expression over \(\alpha\in\Delta_m\) gives the minimum over \(s\) in \((12)\), equivalently the auxiliary inequalities defining \(\gamma\). Letting the separated objective level decrease to the attained primal optimum yields a feasible dual point with the same value. Hence there is no duality gap and \((\mathrm C)\) holds.

At equality, the same calculation yields the checkable complementary conditions
\[
\operatorname{tr}\{X_j^\star F_j(\alpha^\star,t^\star)\}=0,
\qquad
\alpha_s^\star>0
\Longrightarrow
\sum_j\operatorname{tr}(X_j^\star A_{js})=\gamma^\star.
\tag{13}
\]
Therefore any feasible primal--dual pair with equal objectives certifies global optimality.

The exact SDP, trace LP, dual certificate, and law \((\mathrm D)\) constitute the stated \(\mathcal H_{\mathrm{mat}}\) construction. They use \(m=6^q/2\) masses and \(J_{\mathrm{gen}}\) matrix inequalities of order \(n\). Since \(\Pi_n\) already contains every blinded mixture, the implementation constant is one. Matrix-discrepancy or concentration rounding becomes relevant only after adding a product-law, deterministic-law, or externally imposed support restriction.

\textit{Smallest-instance check.}
For \(q=1\) and \(n=4\), one may take
\[
\mathscr S=\{v_1,v_2,v_3\},
\qquad
Q(\alpha)=\sum_{k=1}^3\alpha_kv_kv_k^\top.
\]
Enumeration of the six assignments \(\{\pm v_1,\pm v_2,\pm v_3\}\) gives exactly this triangle of second moments, so \((\mathrm P)\), \((\mathrm T)\), and \((\mathrm D)\) are exact at the smallest population size. In the frozen two-class menu, the product witnesses in \(\texttt{thm:quadruple-menu-frontier}\) make every active score matrix scalar, which is exactly condition \((\mathrm I)\).